\documentclass[12pt]{article}
\usepackage{amsmath}
\usepackage{amsfonts}
\usepackage{amssymb}
\usepackage{geometry}
\usepackage{tikz}
\geometry{a4paper, margin=1in}
\pagenumbering{gobble}


\setlength\parindent{0pt}

\title{02YMECH - Homework 12}
\date{Assigned for the week of Dec 8, 2025}
\author{}
\begin{document}
\maketitle

\section*{Questions}
\begin{enumerate}
    \item Calculate the moment of inertia of a uniform square plate of side length $b$ and mass $M$ about an axis that passes through its center of mass and is perpendicular to the plane of the plate.

    \item Consider a uniform solid cube of side length $b$ and mass $M$.
    \begin{enumerate}
        \item Compute the inertia tensor with respect to axes aligned with the edges of the cube, taking the origin of the coordinate system at the center of mass. Compare your result with the inertia tensor obtained in the lecture for a cube when the origin is chosen at one of its corners.
        \item Since the inertia tensor obtained in part (a) is diagonal in this coordinate system, under what conditions would the angular momentum vector be parallel to the angular velocity vector? Explain briefly.
        \item Interpret the physical meaning of the diagonal elements of the inertia tensor calculated in part (a). Is this result consistent with the moment of inertia of the square plate found in Question~1? Explain the connection.
        \item Suppose the cube is rotating freely in space with no external torques. About which axis (or axes) would you expect the rotation to occur? Based on this, discuss the significance of choosing the center of mass as the origin when calculating the inertia tensor.
    \end{enumerate}
\end{enumerate}


\end{document}
